\scanpage{68}
% archive foliation: 67

$C$ catégorie de Deuring. On a alors la gr-catégorie
$\uline{\operatorname{Eq}}(C)$, qui opère sur $C$, donc
\[
\boxed{\operatorname{Brd}(C)=\pi_0(\uline{\operatorname{Eq}})
\quad\text{opère sur }\deleted{\text{Classes d'isom de }C}
\ \text{Cl. d'isom de }C}
\]
Ceci dit, le quotient est l'ens des classes d'isomorphie des $A$-algèbres $O$ apparentées à $C$
$[$donc, dans le cas des \deleted{\unsure{$A$}} ordres maximaux $\ldots$,
\inserted{le quotient est fini??}
les classes de conjugaison de leurs ordres maximaux,
tandis que $\pi_0(\uline{\operatorname{Eq}})$ correspond \deleted{aux}
\inserted{au groupe des} classes d'isomorphes de $O$-bimodules
$=$ bi-idéaux de $O$ mod aux \unsure{princ.} de \deleted{\unsure{$K$}} $A^*$],
\deleted{\unsure{Noter}} et $\operatorname{Is}(C)$ est l'ens des classes d'isomorphie de $O$-modules à droite ``\unsure{inv.}'' \deleted{\unsure{de}} ou disons des $O$-idéaux \unsure{$*$} à droite $]$.
Dans le cas \inserted{des} $O$ centraux
\deleted{\unsure{sur le}}\inserted{sur $A$}
$[$on dit $C$ centré sur $A$] on a de plus un foncteur \uline{pleinement fidèle}
\diagram{\begin{tikzcd}[row sep=small]
\uuline{\operatorname{Pic}}(A)\arrow[r]&\uline{\operatorname{Eq}}(C)\\
L\arrow[r,mapsto]&(X\longmapsto L\otimes_A X)
\end{tikzcd}}
d'où une suite exacte
\diagram{\begin{tikzcd}[column sep=small]
0\arrow[r,dash]&\operatorname{Pic}(A)\arrow[r]&\operatorname{Brdt}(C)\arrow[r]&
\dfrac{\operatorname{Brdt}(C)}{\operatorname{Pic}(A)}\arrow[r]&0
\end{tikzcd}}
Dans le cas de Brauer (et ordres maximaux $O$)
\[
\frac{\operatorname{Brd}C}{\operatorname{Pic}(A)}
\simeq\frac{\operatorname{Div}(C)}{\operatorname{Div}(A)}
=\coprod_{\substack{p\in\operatorname{Max}A\\p\ \text{ramifié en }C}}
\mathbb Z/e_p\mathbb Z
\]
\margin{$e_p$: ramification en $p$.}
